\section{Safety-Performance Optimization Problem}
	A robotic system is considered ``good'' if it produces high-quality outputs while maintaining safety with high probabilities.
	This paper derived new safety-performance metrics based on Sifat~\etal~\cite{Ash23RAL} to achieve more strict system-level performance measurements and better interpretability, then formulates an online optimization problem to enhance safety and performance.

	\subsection{Safety metric}
    Unlike traditional physical safety concerns (e.g., collision avoidance), this paper focuses on computation safety, which ensures that safety-critical tasks finish their execution before designated deadlines.
    Since this paper mainly focuses on the task scheduling part, we assume all the tasks are designed and implemented correctly.

	\begin{definition}[Safety, \cite{Ash23RAL}]
		A robot system is safe (or schedulable) from the computation perspective if:
		\begin{equation}
			\forall \tau_i, Pr(r_i > D_i) \leq \Theta_i
			\label{eq_safety}
		\end{equation}
		where the threshold $\Theta_i \geq 0$ denotes the probability that the system can tolerate the task $\tau_i$ to miss deadline.
	\end{definition}

	\begin{definition}[Safety Metric]
		A task $\tau_i$'s safety function $\mathcal{S}(\cdot)$ quantifies the penalty (or reward) when the probability of $\tau_i$ missing its deadline $D_i$ rises above (or falls below) $\Theta_i$:
		\begin{equation}
			\mathcal{S}(Pr(r_i > D_i ; \boldsymbol{P}, \boldsymbol{Q}, \textbf{E}) - \Theta_i)
			\label{eq_safety_function}
		\end{equation}
	\end{definition}
    Our experiments adopt the penalty function $\hat{\mathcal{S}}$ of~\cite{Ash23RAL}, but normalize the values into $[0,1]$ (better value interpretability, and enhanced safety-performance combination guarantee, check Section~\ref{section: sp_metric_def}). 
	$\hat{\mathcal{S}}$ has a reward branch when the miss probability is at or below $\Theta_i$ and a steep penalty branch when it exceeds $\Theta_i$, :
	\begin{equation}
		\begin{gathered}
		\mathcal{S}(x) = \frac{\hat{\mathcal{S}}(x) - \hat{\mathcal{S}}(1 - \Theta_i)}{\hat{\mathcal{S}}(0-\Theta_i) - \hat{\mathcal{S}}(1 - \Theta_i)},\\
		\hat{\mathcal{S}}(x) = \begin{cases}
			\log(|x|+1) & x \leq 0 \\
			-0.01\, e^{10|x|}  & x > 0
		\end{cases}
		\end{gathered}
		\label{eq_normalize}
	\end{equation}
	The interpretation formula above guarantees that $\mathcal{S}(x) \in [0,1]$.
    The threshold $\Theta_i$ in \eqref{eq_safety} lets designers tune the tolerable deadline-miss probability per application.

	\subsection{Performance metric}
	A task's performance metric mainly depends on its \emph{QoS budget} (such as an anytime algorithm's running-time limit, which trades execution time for output quality). 

	\begin{definition}[Performance Metric]
    \label{def_perf_metric}
		A computation task $\tau_i$'s performance function $\mathcal{P}(\tau_i, \textbf{E}) \rightarrow [0,1]$ is a normalized function that estimates a task's output quality under an environment described by $\textbf{E}$.
	\end{definition}

	If a task does not have a QoS budget list, then we assume the task's performance metric is always 1.0.

    The performance metric can be measured through offline profiling,
    see Section~\ref{section_generate_tsp_perf} for a concrete example.

	\subsection{SP-Metric}
	\label{section: sp_metric_def}
	The Safety-Performance (SP) metric combines the safety and performance metrics to evaluate the overall performance of a robotic computation system:
	\begin{equation}
		\begin{aligned}
		\textbf{SP}(\boldsymbol{P}, \boldsymbol{Q}; \textbf{E}) \triangleq \sum_{\tau_i \in \boldsymbol{\tau}} & \frac{w_i}{\sum_j w_j} \cdot \mathcal{P}(\tau_i; \textbf{E}) \\
		& \cdot \mathcal{S}\!\left(Pr(r_i > D_i ; \boldsymbol{P}, \boldsymbol{Q}, \textbf{E}) - \Theta_i\right)
		\end{aligned}
		\label{eq_sp_def}
	\end{equation}
	where $w_i > 0$ represents the importance of task $\tau_i$, with the task weights being normalized.
    Both the safety function $\mathcal{S}$ and the performance function $\mathcal{P}$ produce normalized outputs in the range $[0,1]$.

    \begin{Example}
	% TODO, use more meaningful thresholds, such as 0.01 for one task, and 0.6 for another task
        Consider a task set containing two tasks, $\tau_0$ and $\tau_1$, assume the 2 tasks have the following parameters:
    \begin{itemize}
        \item Deadline miss thresholds: $\Theta_0 = 0.5$, $\Theta_1 = 0.5$.
        \item Actual deadline miss ratios: $Pr(r_0 > D_0) = 0.3$, $Pr(r_1 > D_1) = 0.8$.
        \item Performance metrics: $\mathcal{P}(\tau_0) = 0.5$, $\mathcal{P}(\tau_1) = 1.0$.
        \item Task weights: $w_0 = w_1 = 1$.
    \end{itemize}

    Using \eqref{eq_sp_def} with $\Theta_i = 0.5$, the SP-metric is:
    \begin{align}
        & \frac{1}{1+1} \cdot 0.5\cdot \text{Normalize}(\log(1.2)) \ + \nonumber \\
        & \frac{1}{1+1} \cdot 1.0 \cdot \text{Normalize}(-0.01 e^{10 \cdot 0.3})
    \end{align}
    where $\text{Normalize}(\cdot)$ is defined in \eqref{eq_normalize}: for $\tau_0$, $Pr(r_0 > D_0) - \Theta_0 = -0.2 \leq 0$, so $\hat{\mathcal{S}}$ takes the reward branch $\log(1.2)$; for $\tau_1$, $0.3 > 0$, so it takes the penalty branch $-0.01 e^{10 \cdot 0.3}$.


\begin{figure}[htbp]
\centering
\includegraphics[width=0.45\textwidth]{pictures/SP_threshold_plot.pdf}
\caption{SP values with different deadline miss ratios, the deadline miss threshold is 0.5 in the figure.
}
\label{fig_sp_threshold}
\end{figure}
    \end{Example}

    \subsection{Comoparison against the existing SP metric}
    \textbf{Improved Guarantee for Both Safety and Performance.}
    The approach~\cite{Ash23RAL} used an additive combination of safety and performance metrics, which could lead to misleading interpretations.
    High performance values might compensate for low safety values, falsely indicating adequate system behavior; conversely, a trivial algorithm with minimal execution time could easily meet safety requirements yet produce useless output (\eg a SLAM algorithm that always returns the same GPS location runs extremely fast but is wrong).
    The new SP metric addresses these issues by using multiplication to combine the safety and performance metrics.
    Because the per-task term is a product $\mathcal{P} \cdot \mathcal{S}$, a high value on one factor cannot compensate for a low value on the other: for a single task, a per-task term of at least $0.9$ requires both $\mathcal{P}$ and $\mathcal{S}$ to be at least $0.9$ individually, ensuring a balanced trade-off and strong penalization when either value is low.
    We note, however, that this holds per product term; the system-level SP is a \emph{weighted sum} of such terms, so a system-level SP of $0.9$ does not by itself imply that every task's term reaches $0.9$, as a strong task can mask a weak one.


    \textbf{Enhanced interpretability.}
    The previous SP metric~\cite{Ash23RAL} had unbounded values (\eg from $-2000$ to $2$), making it difficult to intuitively interpret the meanings of SP values (such as when the system always meets/misses deadlines) without reference points.
    The new metric normalizes values to $[0,1]$, inherently encoding worst- and best-case scenarios.
    This standardization simplifies threshold setting for robotic engineers.
    The safety term can also be interpreted as a discount to the performance term in the objective function~\eqref{eq_overall_obj}: a higher deadline missing rate would reduce the system's overall performance.

	\subsection{SP Metric Optimization Problem}
    This paper addresses the problem of optimizing the scheduling algorithms of a robotic computation system operating under unknown and dynamic environments.
    Task computation requirements often vary across different environments, necessitating dynamic adjustments to scheduling algorithms to improve the overall safety and performance of the system.
    The optimization problem is defined as follows:

	\begin{equation}
		\begin{aligned}
		\max_{\boldsymbol{P}, \boldsymbol{Q}} \sum_{\tau_i \in \boldsymbol{\tau}} & \frac{w_i}{\sum_j w_j} \cdot \mathcal{P}(\tau_i; \textbf{E}) \\
		& \cdot \mathcal{S}\!\left(Pr(r_i > D_i ; \boldsymbol{P}, \boldsymbol{Q}, \textbf{E}) - \Theta_i\right)
		\end{aligned}
		\label{eq_overall_obj}
	\end{equation}
	\begin{equation}
		\rtDist{i} = \mathcal{R}\!\left(\extDist{i}(\textbf{E};q_i),\; \text{hp}(i;\boldsymbol{\tau},\boldsymbol{P},\boldsymbol{Q})\right)
		\label{eq_prob_rta_in_opt}
	\end{equation}
	\begin{equation}
		Pr(r_i > D_i) \leq \Theta_i, \quad \forall \tau_i \in \boldsymbol{\tau}^{safe}
		\label{eq_important_task_constraint}
	\end{equation}

        The objective \eqref{eq_overall_obj} maximizes the SP metric \eqref{eq_sp_def} over the priority assignment $\boldsymbol{P}$ (under fixed-task priority scheduling) and the QoS budgets $\boldsymbol{Q}$ (here, the running-time limits of anytime algorithms), given the environment vector $\textbf{E}$.
		Equation~\eqref{eq_prob_rta_in_opt} computes the response-time distribution $\rtDist{i}$ from $\tau_i$'s execution-time distribution $\extDist{i}(\textbf{E};q_i)$ and its higher-priority task set.
        $\text{hp}(i;\boldsymbol{\tau},\boldsymbol{P},\boldsymbol{Q})$ denotes the higher-priority set $\text{hp}(i)$ of \eqref{eq: rta_scalar} in the task set $\boldsymbol{\tau}$.

        Not all tasks are equally safety-critical.
        We therefore designate a subset $\boldsymbol{\tau}^{safe} \subseteq \boldsymbol{\tau}$ of \emph{important tasks}: the tasks whose safety thresholds must be honored as hard constraints.
        Concretely, $\boldsymbol{\tau}^{safe}$ is taken as the tasks carrying the highest safety weights, reflecting that the tasks the designer cares most about should not have their deadlines violated.
        Equation~\eqref{eq_important_task_constraint} imposes this as an explicit constraint of the optimization: for every important task $\tau_i \in \boldsymbol{\tau}^{safe}$, the deadline-miss probability under the chosen $\{\boldsymbol{P}, \boldsymbol{Q}\}$ must not exceed its safety threshold $\Theta_i$ (the same $\Theta_i$ defined in \eqref{eq_safety}).
        The non-safety-critical tasks contribute to the system performance objective \eqref{eq_overall_obj} but are not individually constrained, so the scheduler may trade their performance against overall safety-performance when necessary.

		% This overview is not ncessary in problem definition
        % The optimization in \eqref{eq_overall_obj} is solved online by a lightweight incremental algorithm that interleaves priority assignment and QoS-budget optimization; its design is presented in Sections~\ref{section_priority_opt} and~\ref{section_config_opt}.
        % The constraint in \eqref{eq_important_task_constraint} is \emph{not} left to the objective to satisfy softly: the search is constrained to maintain a certified safe fallback so that, at every shipped interval, each important task satisfies $Pr(r_i > D_i) \leq \Theta_i$.
        % This important-task safety guarantee is a built-in property of the framework---not a conditional assumption on a user-supplied execution-time upper bound---and is detailed in Section~\ref{section_safety_fallback}.





	\subsection{Challenges of SP-metric optimization}

    The optimization problem in \eqref{eq_overall_obj} presents several challenges:
\begin{itemize}
    \item \textbf{Mixed combinatorial and continuous variables:} The problem involves both continuous variables (QoS budgets $\boldsymbol{Q}$) and combinatorial variables (the priority assignment $\boldsymbol{P}$).
    \item \textbf{Complicated objective function:} The SP metric's formulation \eqref{eq_sp_def} includes probabilistic components and lacks properties like convexity or continuity, complicating the optimization process.
    \item \textbf{Limited computation resources:} \eqref{eq_overall_obj} must be solved online under unknown environments, yet mobile robots carry limited computation resources largely consumed by functional tasks such as motion planning and SLAM, so the solver must be lightweight---ideally linear in run-time complexity.
\end{itemize}
    These properties rule out general-purpose frameworks such as convex optimization or integer linear programming, as well as the exponential-complexity offline methods of~\cite{zhao2018unified, Wang2023OptimizingLET, Verucchi2020LatencyAwareGO}.
    A customized online algorithm is therefore required.

	
	\subsection{Solution Framework Overview}
    \label{section_solution_overview}
    The online problem~\eqref{eq_overall_obj} is addressed by a lightweight framework built on two ideas.
    First, since $\textbf{E}$ changes continuously and re-optimization intervals are short, only a few tasks' distributions change slowly between invocations, so we can each warm-starts an incremental search from the last deployed configuration, and perform incremental optimization to focus on what's changed between scheduelr invocation.
    Second, the hard constraint~\eqref{eq_important_task_constraint} is enforced by a certified safe fallback---a configuration verified offline via worst-case response-time analysis to satisfy $Pr(r_i > D_i) \leq \Theta_i$ for every $\tau_i \in \boldsymbol{\tau}^{safe}$ and retained across invocations;
	a candidate is adopted only if it strictly improves the objective \emph{and} satisfies the constraint, otherwise the fallback is shipped.
    Sections~\ref{section_priority_opt}--\ref{section_safety_fallback} develop the priority, QoS, and fallback components.
